\documentclass[conference]{IEEEtran}
\IEEEoverridecommandlockouts
\usepackage{cite}
\usepackage{amsmath,amssymb,amsfonts}
\usepackage{algorithmic}
\usepackage{graphicx}
\usepackage{textcomp}
\usepackage{xcolor}
\usepackage{booktabs}
\usepackage{multirow}
\def\BibTeX{{\rm B\kern-.05em{\sc i\kern-.025em b}\kern-.08em
    T\kern-.1667em\lower.7ex\hbox{E}\kern-.125emX}}
\begin{document}

\title{Perceived Stress Detection Using Wearable
Physiological Sensors with Personalized Adaptive
Machine Learning}

\author{
\IEEEauthorblockN{1\textsuperscript{st} Wehara Soyza}
\IEEEauthorblockA{
\textit{Department of Information Technology} \\
\textit{Sri Lanka Institute of Information Technology}\\
Malabe, Sri Lanka \\
wehara.soyza@email.com}
\and
\IEEEauthorblockN{2\textsuperscript{nd} Supervisor Name}
\IEEEauthorblockA{
\textit{Department of Information Technology} \\
\textit{Sri Lanka Institute of Information Technology}\\
Malabe, Sri Lanka \\
supervisor@sliit.lk}
}

\maketitle

% ── ABSTRACT ─────────────────────────────────────────────
\begin{abstract}
Chronic and undetected stress is a leading contributor to
cardiovascular, metabolic, and mental health disorders, yet
most wearable stress-detection systems only report a
subject's state after stress has already fully manifested,
offering the user no opportunity to intervene earlier. This
work presents a wearable physiological stress-monitoring
system that addresses three compounding gaps in the
literature: the absence of an early-warning signal before
stress onset, the reliance on population-level thresholds
that ignore individual physiological variability, and the
lack of adaptation as a user's baseline naturally drifts
over time. Using the WESAD dataset (15 subjects, 700\,Hz
chest RespiBAN recordings), a binary Normal/Stress label
aligned with the Perceived Stress Scale (PSS) is constructed
directly from the dataset's baseline and TSST-stress
conditions, and a 22-feature representation spanning heart
rate variability, electrodermal activity, and motion is
selected via a four-method statistical vote. An XGBoost
classifier evaluated with leave-one-subject-out
cross-validation forms the static baseline (F1-macro 0.617),
on top of which a personalization framework is introduced:
per-user resting-baseline calibration combined with an
exponential-moving-average continual-learning layer that
adapts the baseline session-by-session using only
confidently-Normal windows, supplemented by a lightweight
trend-monitoring indicator that flags a rising physiological
trajectory ahead of the binary label itself changing.
Personalization lifts F1-macro to 0.628, Ordinal Mean Absolute
Error from 0.357 to 0.194, and overall accuracy by 6 percentage
points, with the largest individual gain on the Stress class
(F1-Stress +11.3\%). A cross-dataset check on SWELL
(HRV-only, no electrodermal signal) confirms that
electrodermal activity is not a redundant input but an
essential one for reliable detection. These results and the
accompanying personalization framework are packaged into a
low-cost IoT smart ring (XIAO ESP32-C3, MAX30102, Grove GSR,
MPU-6050) intended to bring personalized, adaptive stress
detection to continuous, real-world wear rather than
laboratory conditions alone.
\end{abstract}

\begin{IEEEkeywords}
stress detection, personalized machine learning,
wearable sensors, continual learning,
electrodermal activity, heart rate variability,
stress trend monitoring, smart ring, IoT health monitoring
\end{IEEEkeywords}

% ── SECTION 1 — INTRODUCTION ─────────────────────────────
\section{Introduction}

Stress has become one of the most pervasive and costly
public health burdens of modern life. Prolonged exposure to
unmanaged stress is a well-established contributor to
cardiovascular disease, weakened immune function, sleep
disruption, and the onset or worsening of anxiety and
depressive disorders. Because stress accumulates gradually
and its physiological consequences are frequently invisible
to the person experiencing it until symptoms are already
severe, there is a clear and growing need for tools that can
surface a user's stress state continuously, unobtrusively,
and, ideally, before it has fully developed rather than only
after the fact.

Wearable physiological sensors have emerged as the most
practical route to this kind of continuous monitoring.
Heart rate variability (HRV), derived from photoplethysmography
(PPG) or electrocardiography (ECG), and electrodermal
activity (EDA), measured via galvanic skin response (GSR),
are the two signal families most consistently linked to
autonomic nervous system arousal, and both can now be
captured from small, low-power, comfortably-worn devices.
A substantial body of recent work has demonstrated that
machine learning models trained on HRV and EDA features can
classify a wearer's stress state with reasonable accuracy,
and this combination has become the dominant sensing
approach in the field.

Despite this progress, existing systems share a central
limitation concerning individual variability. Physiological
baselines (resting heart rate, resting skin conductance, and
the sensitivity of each to arousal) vary substantially between
individuals, and even within the same individual over time.
Systems trained and evaluated purely on population-level
thresholds implicitly assume a "typical" physiology that no
real user exactly matches, and these systems are, without
exception in the literature reviewed here, static after
training: none adapt to a specific wearer, and none continue
to adapt as that wearer's own baseline naturally drifts across
weeks or months of continued use. A further, secondary
limitation is that these systems report stress only as a
discrete, present-tense classification, with no mechanism for
surfacing a rising trend ahead of that classification itself
changing. This is a minor extension addressed briefly in this
work (Section~\ref{sec_trend}), subordinate to the
personalization gap above, which is this paper's central
focus.

This work is structured around three research objectives that
directly target these gaps: (1) establish a static binary
Normal/Stress classifier as a rigorously-validated baseline,
using leave-one-subject-out cross-validation rather than the
less conservative k-fold splits common in prior work; (2)
build a personalized, continually-adaptive version of that
model, the central contribution of this work, comprising
per-user baseline calibration and an exponential-moving-average
continual-learning layer, supplemented by a lightweight
trend-monitoring indicator layered on top; and (3) validate
the complete pipeline on a purpose-built, low-cost IoT smart
ring, moving the system from laboratory data toward
real-world, continuous wear.

The remainder of this paper is organized as follows. Section
II reviews the related literature and distills the specific
gaps this work addresses. Section III details the dataset,
preprocessing pipeline, feature selection, and the methodology
for each of the three components above, including the smart
ring hardware. Section IV reports results and validation:
WESAD-based evaluation, IoT ring hardware validation, and
consistency with existing research. Section V discusses the
implications of
these results, compares this work against the closest prior
systems, and states its limitations. Section VI summarizes
the paper's novel contributions against the literature gaps
identified in Section II, and Section VII concludes with a
summary of findings and directions for future work.

% ── SECTION 2 — RELATED WORK ─────────────────────────────
\section{Related Work}

\subsection{Physiological Stress Detection Using
Wearable Sensors}
Wearable-sensor-based stress detection has matured into an
active research area, with HRV and GSR emerging as the
dominant combination of input signals across the studies
reviewed here \cite{b1}\cite{b3}\cite{b4}\cite{b14}. The
large majority of these systems frame stress detection as a
binary classification problem (Normal versus Stress) using
supervised learning over hand-crafted or lightly-engineered
features. While reported accuracies are frequently strong on
held-out data from the same population, none of these systems
attempt to detect the onset of stress before it is fully
established; each reports the wearer's state only after the
underlying model has already classified the current window as
Stress, offering no earlier signal a user could act on.

\subsection{Electrodermal Activity and HRV
Feature Extraction}
A consistent set of features recurs across the literature as
the most discriminative for stress: tonic EDA measures such
as skin conductance level (SCL), phasic measures such as skin
conductance response (SCR) count and amplitude, and HRV
time-domain statistics such as average NN interval (AVNN),
standard deviation of heart rate (SDHR), and the percentage of
successive NN interval differences greater than 20\,ms
(pNN20) \cite{b2}\cite{b5}\cite{b7}\cite{b12}. Preprocessing
pipelines for isolating these features from raw PPG/ECG and
GSR signals are well validated in this literature. However,
the feature values and the decision thresholds applied to
them are, almost universally, computed and validated at the
population level; no per-subject recalibration is performed,
which is precisely the limitation the personalization layer
in this work is designed to address.

\subsection{Personalized Stress Models}
A smaller subset of the literature explicitly acknowledges
that individual physiological variability limits the
reliability of population-trained models
\cite{b6}\cite{b8}\cite{b9}\cite{b11}, and several of these
studies argue for some form of per-subject baseline as a
necessary correction. Even where this is proposed, the
resulting personalized model is fixed once calibrated: none
of the reviewed systems revisit or update that personal
baseline after its initial establishment, meaning any
subsequent drift in the wearer's physiology (from changes in
fitness, medication, sleep patterns, or simply the passage of
time) is never reflected in the deployed model.

\subsection{Continual Learning and Adaptive Systems}
Continual or lifelong learning, in which a deployed model
continues to update itself using new data collected during
normal operation, is well studied in machine learning
generally, and is referenced as a desirable property in
several stress-detection papers \cite{b3}\cite{b4}\cite{b6}.
None of the systems reviewed for this work, however, actually
implement such a mechanism for stress detection specifically.
This represents a complete gap in the applied literature: the
idea is acknowledged as valuable, but no existing wearable
stress system has published an implementation, evaluation, or
even a concrete lightweight mechanism (such as an
exponential moving average) suited to the resource constraints
of a wearable device.

\subsection{IoT Wearable Systems for Stress Monitoring}
A parallel thread of work focuses on the systems-engineering
side of wearable stress detection: low-power microcontrollers,
edge or near-edge inference, and real-time data pipelines from
sensor to prediction \cite{b7}\cite{b13}\cite{b14}. Of
particular relevance, one study validates a hardware
configuration closely related to the one adopted in this work
and reports 91.3\% accuracy \cite{b10}, confirming the
practical viability of this sensor combination on
resource-constrained hardware. However, these systems focus
purely on the engineering pipeline; none incorporate
personalization or the trend-detection capability introduced
here, meaning the hardware validation and the modeling
contributions in the literature have so far remained
disconnected from one another.

\subsection{Research Gaps}
Table~\ref{tab_gaps} summarizes the six gaps distilled from
the fifteen studies reviewed above. Gaps G1 and G3, the
absence of any early trend warning and the fact that static
models are never revisited after training, are present in
every study surveyed, without exception. G2 and G4 concern
individual variability and motion-induced signal confounds,
both acknowledged in a subset of the literature but rarely
addressed end-to-end. G5 and G6 are methodological: most
studies construct labels from discrete experimental events
rather than aligning them to a validated self-report
instrument such as the Perceived Stress Scale (PSS), and the large
majority evaluate using k-fold cross-validation, which,
because folds mix windows from the same subject across
training and test splits, tends to inflate reported accuracy
relative to the stricter leave-one-subject-out (LOSO)
protocol used in this work.

\begin{table}[htbp]
\caption{Key Research Gaps Identified
Across 15 Studies}
\begin{center}
\begin{tabular}{|c|p{3.5cm}|p{2.5cm}|}
\hline
\textbf{Gap} &
\textbf{Description} &
\textbf{Source Papers} \\
\hline
G1 & No early stress trend warning &
All 15 papers \\
G2 & Universal models ignore individual
baseline variability &
\cite{b1}\cite{b4}\cite{b8}\cite{b10} \\
G3 & Static models degrade over time &
All 15 papers \\
G4 & Physical activity confounds signals &
\cite{b1}\cite{b4}\cite{b5}\cite{b8} \\
G5 & Event-based labels not PSS-aligned &
\cite{b1}\cite{b4}\cite{b5}\cite{b9} \\
G6 & K-fold inflation, LOSO rarely used &
\cite{b7}\cite{b15} \\
\hline
\end{tabular}
\label{tab_gaps}
\end{center}
\end{table}

% ── SECTION 3 — METHODOLOGY ──────────────────────────────
\section{Methodology}

\subsection{Dataset and Label Creation}
This work uses the WESAD dataset, comprising synchronized
chest-worn RespiBAN recordings (ECG, EDA, and three-axis
accelerometry, all at 700\,Hz) from 15 subjects across a
resting baseline condition and a Trier Social Stress Test
(TSST) condition. Rather than treating WESAD's original
multi-class protocol labels as the target directly, a binary
label aligned with the intent of the Perceived Stress Scale
is constructed: the baseline condition is mapped to Normal
(0) and the TSST stress condition is mapped to Stress (1). A
30-second exclusion window is applied around every true
condition transition to remove ambiguous boundary samples,
and any candidate window whose underlying samples span more
than one condition is discarded outright rather than
retained with a averaged or majority label. This yields 854
valid 60-second windows (30-second step, i.e. 50\% overlap
between consecutive windows), summarized in
Table~\ref{tab_dataset}: 587 Normal (68.7\%) and 267 Stress
(31.3\%).

\begin{table}[htbp]
\caption{WESAD Dataset Window Statistics}
\begin{center}
\begin{tabular}{|l|c|c|}
\hline
\textbf{Class} &
\textbf{Windows} &
\textbf{Percentage} \\
\hline
Normal (0) & 587 & 68.7\% \\
Stress (1) & 267 & 31.3\% \\
\hline
\textbf{Total} & \textbf{854} & \textbf{100\%} \\
\hline
\multicolumn{3}{l}{Window: 60s, Step: 30s,
±30s boundary excluded} \\
\end{tabular}
\label{tab_dataset}
\end{center}
\end{table}

\subsection{Signal Preprocessing Pipeline}

\subsubsection{ECG Preprocessing and HRV Extraction}
The raw chest ECG signal is bandpass filtered (0.5–5\,Hz) and
DC-corrected using a 2-second moving average to remove slow
baseline wander. R-peaks are detected using an adaptive,
six-stage fallback procedure, beginning with a standard
Pan-Tompkins-style detector and progressively relaxing
detection thresholds (and, as a last resort, testing the
inverted signal) for subjects whose ECG morphology defeats the
primary detector, since a small number of subjects in WESAD
produce atypical waveforms that a single fixed-threshold
detector fails on entirely. The resulting inter-beat interval
(IBI) sequence is constrained to a hard physiological range of
300–2000\,ms (200 to 30 beats per minute), and a
median-based contextual outlier filter rejects any IBI value
deviating more than 20\% from its local 5-beat median; rejected
samples are replaced via linear interpolation rather than
removed outright, to preserve a continuous time series for
downstream feature computation.

\subsubsection{EDA Preprocessing}
The raw EDA signal undergoes interquartile-range (IQR) outlier
removal, followed by linear detrending to correct for slow
electrode drift. A Butterworth low-pass filter (5\,Hz cutoff)
isolates the smoothed signal, and a subsequent high-pass stage
(0.05\,Hz) separates the tonic skin conductance level (SCL)
component from the faster phasic skin conductance response
(SCR) component. First- and second-order derivatives of the
filtered signal are computed to capture the rate and
acceleration of change in conductance, and a 5-level Haar
wavelet decomposition provides multi-resolution energy
features sensitive to both slow tonic drift and rapid phasic
bursts.

\subsubsection{Accelerometer Motion Gating}
Tri-axial accelerometer samples are reduced to a single
resultant magnitude
\begin{equation}
M = \sqrt{a_x^2 + a_y^2 + a_z^2}
\label{eq_acc}
\end{equation}
which serves two distinct purposes in this pipeline. First, any
window in which more than a set proportion of samples exceed
a 1.2\,g threshold is flagged as motion-contaminated and
excluded from training, protecting the HRV and EDA features
from gross-motion artifacts. Second, the same accelerometer
statistics (mean and standard deviation of $M$) are themselves
retained as predictive features rather than discarded once
gating is complete, since subtle motion level (restlessness,
fidgeting) carries independent stress-relevant information
beyond simply flagging unusable windows. Because WESAD's
chest-mounted sensor placement is comparatively stable during
both the baseline and TSST conditions, motion contamination
in this dataset was confirmed at 0\%, meaning the gating logic
was exercised without discarding any WESAD windows on this
basis; the same logic is expected to activate more frequently
on the wrist/finger-worn deployment hardware described in
Section~\ref{sec_hardware}.

\subsection{Feature Extraction and Selection}
Forty candidate features are computed across the HRV, EDA,
and accelerometer domains described above. Feature selection
is performed via a four-method statistical vote: (1) Spearman
rank correlation between each feature and the binary label,
(2) Extra Trees ensemble feature importance, (3) recursive
feature elimination (RFE) using a Random Forest estimator, and
(4) permutation importance. A candidate feature is retained in
the final set only if it receives votes from at least two of
the four methods, yielding the 22-feature representation
summarized in Table~\ref{tab_features} and used for every
model reported in Section IV.

\begin{table}[htbp]
\caption{Final 22 Selected Features by Category}
\begin{center}
\begin{tabular}{|p{1.3cm}|p{5.1cm}|c|}
\hline
\textbf{Category} & \textbf{Features} & \textbf{Count} \\
\hline
HRV & AVNN, MeanHR, SDHR, pNN20, SD1 & 5 \\
\hline
EDA Tonic & SCL\_mean, EDA\_mean, EDA\_slope, EDA\_std & 4 \\
\hline
EDA Phasic & SCR\_count, SCR\_amp\_mean, SCR\_width\_mean & 3 \\
\hline
EDA Stat. & EDA\_q10, EDA\_q75, EDA\_q90, EDA\_range & 4 \\
\hline
EDA Deriv. & deriv1\_std, deriv2\_mean & 2 \\
\hline
EDA Wavelet & wavelet\_d1\_energy, wavelet\_d3\_energy & 2 \\
\hline
ACC & ACC\_mean, ACC\_std & 2 \\
\hline
\textbf{Total} & & \textbf{22} \\
\hline
\end{tabular}
\label{tab_features}
\end{center}
\end{table}

\subsection{Static Binary Classification}
The static baseline model is evaluated using
leave-one-subject-out (LOSO) cross-validation across all 15
WESAD subjects: for each fold, a model is trained on 14
subjects and evaluated on the fully held-out 15th, so that no
window from the test subject ever influences training in any
form, avoiding the accuracy inflation associated with
subject-mixed k-fold splits (Gap G6). Class imbalance is
handled via class weighting during training rather than
resampling, and feature standardization
(\texttt{StandardScaler}) is fit exclusively on each fold's
training subjects and then applied to the held-out subject,
preventing any information leakage from the test subject into
the scaling statistics. Ten candidate classifiers are compared
under this identical protocol; XGBoost is selected as the
deployed static model on the basis of the highest F1-macro
score among the ten (Section~\ref{sec_results_obj1}).

\subsection{Personalized Adaptive Model}

\subsubsection{Layer 1: Personalized Baseline
Calibration}
Population-level feature thresholds fail to account for the
substantial resting-state variability between individuals:
two people can have meaningfully different resting heart
rates or skin conductance levels while both being equally
"Normal" for themselves. To correct for this, a short
(approximately 2-minute) resting calibration period is
introduced: the first three windows a user's model classifies
as Normal are averaged to establish a personal mean and
standard deviation for every feature. Subsequent readings are
normalized against this personal baseline rather than against
a population statistic. EDA-domain features are normalized as
a ratio to the personal mean (since their absolute scale is
strongly influenced by individual skin and electrode
characteristics), while HRV and accelerometer features are
z-score normalized:
\begin{equation}
z_{HRV/ACC} = \frac{x - \mu_{personal}}
{\sigma_{personal}}
\label{eq_zscore}
\end{equation}
\begin{equation}
r_{EDA} = \frac{x}{\mu_{personal}}
\label{eq_ratio}
\end{equation}

\subsubsection{Layer 2: EMA Continual Learning}
A one-time calibration is itself insufficient, since an
individual's physiological baseline is not fixed; it drifts
gradually with changes in fitness, sleep, medication, or
simply the passage of time. To track this drift without
requiring the user to ever recalibrate manually, the personal
baseline is updated after every session using an exponential
moving average:
\begin{equation}
\hat{\mu}_{new} = \alpha \cdot
\bar{x}_{session} +
(1-\alpha) \cdot \hat{\mu}_{old},
\quad \alpha = 0.2
\label{eq_ema}
\end{equation}
Critically, this update is computed only from windows the
model itself classified as Normal within that session (with a
minimum of three such windows required); sessions containing
fewer than three confidently-Normal windows are skipped
entirely, so that a genuinely stressful session cannot itself
contaminate the resting baseline it is being compared against.
With $\alpha = 0.2$, the baseline adapts meaningfully over
roughly 10–15 sessions, corresponding to approximately three
weeks of regular use, fast enough to track real physiological
drift but slow enough that no single anomalous session can
distort the baseline.

\subsubsection{Layer 3: Trend Monitoring}
\label{sec_trend}
As a lightweight supplement to the two personalization layers
above, which remain this work's central contribution, a
rolling 5-window feature-direction vote monitors whether a
majority of tracked HRV/EDA/ACC features are moving in a
stress-consistent direction:
\begin{equation}
V_{trend} = \frac{f_{stress\_direction}}
{f_{total}} \geq \theta_{trend},
\quad \theta_{trend} = 0.35
\label{eq_trend}
\end{equation}
When this vote fires, a separate Stress-Trending flag is
raised alongside (never in place of) the binary
Normal/Stress label, since no ground-truth "pre-stress" class
exists in WESAD (or in any validated self-report instrument)
against which such an override could itself be validated. Its
role is intentionally minor: an advisory early signal, not an
alteration of the validated binary classification.

\subsection{IoT Smart Ring Hardware}
\label{sec_hardware}
The complete pipeline is deployed on a purpose-built,
low-cost IoT ring, summarized in
Table~\ref{tab_hardware}. A Seeed XIAO ESP32-C3 serves as the
onboard microcontroller and WiFi radio; a MAX30102
reflective-PPG sensor (I2C address \texttt{0x57}) captures the
pulse signal from which HRV features are derived; a Grove GSR
module (read via a single analog input) captures electrodermal
activity from a two-electrode finger circuit; and an MPU-6050
six-axis IMU (I2C address \texttt{0x68}, with its address-select
pin tied to ground) provides the accelerometer data used for
both motion gating and the ACC\_mean/ACC\_std features. All
three sensors, along with the microcontroller, are powered from
a single regulated 3.3\,V rail sourced from a 250\,mAh LiPo
battery managed by a TP4056 charge-protection module, ensuring
consistent sensor operation whether the device is
battery-powered or connected via USB. Raw sensor windows are
streamed over WiFi/MQTT to a Flask inference backend, which
performs feature extraction identical to the offline WESAD
pipeline described above before returning a prediction.

\begin{table}[htbp]
\caption{MindPulse Smart Ring Hardware Components}
\begin{center}
\begin{tabular}{|l|l|l|}
\hline
\textbf{Component} &
\textbf{Sensor} &
\textbf{Role} \\
\hline
XIAO ESP32-C3 & MCU + WiFi &
Processing, MQTT TX \\
MAX30102 & PPG sensor &
HRV (I2C 0x57) \\
Grove GSR & EDA sensor &
Skin conductance \\
MPU-6050 & Accelerometer &
Motion gating (I2C 0x68) \\
LiPo 250mAh & Battery &
Wireless power \\
TP4056 & Charge module &
LiPo protection \\
\hline
\end{tabular}
\label{tab_hardware}
\end{center}
\end{table}

% ── SECTION 4 — RESULTS AND VALIDATION ───────────────────
\section{Results and Validation}

\subsection{Static Model Comparison}
\label{sec_results_obj1}
Table~\ref{tab_models} reports LOSO cross-validation results
for all ten candidate classifiers. While several models,
notably MLP (accuracy 0.883) and LDA (accuracy 0.838), achieve
higher raw accuracy than XGBoost, this is substantially driven
by strong performance on the majority Normal class rather than
balanced performance across both classes; XGBoost achieves the
highest F1-macro score (0.617) among all ten models, reflecting
the most balanced trade-off between Normal and Stress
detection, and is therefore selected as the deployed static
model on this basis rather than on raw accuracy alone.

\begin{table}[htbp]
\caption{LOSO Cross-Validation:
10 Model Comparison
(854 windows, 15 WESAD subjects)}
\begin{center}
\begin{tabular}{|l|c|c|c|c|}
\hline
\textbf{Model} &
\textbf{Acc} &
\textbf{F1-Mac} &
\textbf{F1-Nor} &
\textbf{F1-Str} \\
\hline
\textbf{XGBoost} &
\textbf{0.796} &
\textbf{0.617} &
\textbf{0.877} &
\textbf{0.689} \\
SVM-Linear & 0.728 & 0.544 & 0.878 & 0.550 \\
LDA & 0.838 & 0.615 & 0.935 & 0.708 \\
Logistic Reg & 0.770 & 0.579 & 0.906 & 0.628 \\
SVM-RBF & 0.790 & 0.573 & 0.886 & 0.674 \\
Grad. Boost & 0.796 & 0.577 & 0.883 & 0.702 \\
AdaBoost & 0.757 & 0.539 & 0.855 & 0.642 \\
Rand. Forest & 0.830 & 0.585 & 0.894 & 0.748 \\
MLP & 0.883 & 0.583 & 0.936 & 0.813 \\
KNN & 0.823 & 0.544 & 0.891 & 0.740 \\
\hline
\multicolumn{5}{l}{Bold = selected (highest F1-macro)} \\
\end{tabular}
\label{tab_models}
\end{center}
\end{table}

\subsection{Confusion Matrix}
Table~\ref{tab_cm} presents the pooled LOSO confusion matrix
for the selected XGBoost model. Of 587 true Normal windows,
500 (85.2\%) are correctly classified; of 267 true Stress
windows, 194 (72.7\%) are correctly classified. The higher
false-negative rate on the Stress class relative to Normal is
consistent with the class imbalance in the underlying dataset
(Table~\ref{tab_dataset}) and motivates the personalization
results reported next.

\begin{table}[htbp]
\caption{XGBoost Confusion Matrix:
LOSO Evaluation}
\begin{center}
\begin{tabular}{|l|c|c|}
\hline
\multirow{2}{*}{\textbf{Actual}} &
\multicolumn{2}{c|}{\textbf{Predicted}} \\
\cline{2-3}
& \textbf{Normal} &
\textbf{Stress} \\
\hline
Normal & 500 & 87 \\
Stress & 73 & 194 \\
\hline
\end{tabular}
\label{tab_cm}
\end{center}
\end{table}

\subsection{Personalized Model Results}
Table~\ref{tab_comparison} compares the static model
against the fully personalized pipeline
(baseline calibration plus EMA continual learning) under the
same LOSO protocol. Personalization improves every reported
metric: accuracy rises 6.0 percentage points (79.6\% to
85.6\%), F1-macro improves from 0.617 to 0.628, and Ordinal
Mean Absolute Error (OrdMAE) falls from 0.357 to 0.194. OrdMAE
is the average absolute deviation between predicted and true
class index, and it penalizes confident misclassifications
more than borderline ones. The largest single-class improvement is
on Stress detection specifically (F1-Stress +0.113, i.e.
+11.3 percentage points), which is precisely the class the
static model struggled with most in Table~\ref{tab_cm},
indicating that individual baseline normalization
disproportionately helps distinguish a user's genuine stress
response from their own resting variability.

\begin{table}[htbp]
\caption{Static vs Personalized
Performance Comparison}
\begin{center}
\begin{tabular}{|l|c|c|c|}
\hline
\textbf{Metric} &
\textbf{Static} &
\textbf{Personal} &
\textbf{Change} \\
\hline
Accuracy & 79.6\% & 85.6\% & +6.0\% ✓ \\
F1-Macro & 0.617 & 0.628 & +0.011 ✓ \\
F1-Normal & 0.877 & 0.938 & +0.061 ✓ \\
F1-Stress & 0.689 & 0.802 & +0.113 ✓ \\
OrdMAE & 0.357 & 0.194 & -0.163 ✓ \\
\hline
\end{tabular}
\label{tab_comparison}
\end{center}
\end{table}

\subsection{Trend Monitoring: Preliminary Note}
As a supplementary indicator rather than a focus of this
evaluation, the trend-monitoring layer (Eq.~\ref{eq_trend})
has so far been validated only on controlled synthetic
scenarios, correctly remaining silent on stable and on
noisy/undirected signals and firing only once a consistent
multi-feature rise is sustained across the 5-window buffer.
Session-level effectiveness on real ring data is left to
future work (Section~VII).

\subsection{Cross-Dataset Validation (SWELL)}
To test whether electrodermal activity is a genuinely
necessary input rather than a convenient but replaceable one,
the WESAD-trained model is re-evaluated using only the five
HRV features common to WESAD and the independent SWELL
knowledge-work dataset, which contains no EDA signal at all.
Table~\ref{tab_swell} shows a substantial performance
collapse under this restriction: F1-macro falls from 0.617
(full 22-feature WESAD model) to 0.270, and F1-Stress falls to
0.165, only marginally better than chance. This confirms that
electrodermal activity is not a redundant or substitutable
input in this pipeline, and directly justifies the ring's
finger-worn form factor (Section~\ref{sec_hardware}), since a
wrist-only device with HRV sensing alone would be expected to
perform little better than the degraded HRV-only result shown
here.

\begin{table}[htbp]
\caption{Cross-Dataset Validation on SWELL
(5 HRV Features, No GSR Available)}
\begin{center}
\begin{tabular}{|l|c|c|}
\hline
\textbf{Metric} &
\textbf{WESAD LOSO} &
\textbf{SWELL Cross} \\
\hline
F1-Macro & 0.617 & 0.270 \\
F1-Normal & 0.877 & 0.593 \\
F1-Stress & 0.689 & 0.165 \\
OrdMAE & 0.357 & 0.925 \\
\hline
\multicolumn{3}{l}{Confirms GSR essential
for reliable stress detection} \\
\end{tabular}
\label{tab_swell}
\end{center}
\end{table}

\subsection{IoT Ring Hardware Validation}
The firmware and software pipeline (sensor sampling, MQTT
streaming, Flask-side feature extraction, and prediction) has
been confirmed operational end-to-end on the physical device.
Full end-to-end system characterization on the assembled ring,
sensor-to-prediction latency, MAX30102-derived heart rate
accuracy against a reference pulse oximeter, Grove GSR
responsiveness, MPU-6050 motion-gate accuracy, and measured
battery life: [TO BE FILLED].

\subsection{Validation Against Existing Research}
The results above are broadly consistent with the literature
reviewed in Section II. The 79.6\% static accuracy and 85.6\%
personalized accuracy reported here are comparable to the
91.3\% accuracy reported by the most hardware-similar prior
system \cite{b10}, which uses the same core sensor combination
(ESP32, MAX30102, Grove GSR) but a dual-branch CNN classifier
and no personalization layer; the remaining gap is consistent
with that system's use of a more complex, higher-capacity
model rather than any inconsistency in this work's own
pipeline. The feature domains selected in this work, HRV and
EDA, match those identified as most discriminative across the
broader literature \cite{b1}\cite{b2}\cite{b5}\cite{b7}\cite{b12},
supporting the validity of the four-method feature-selection
procedure in Section III. Finally, the accuracy improvement
observed from personalization is directionally consistent
with the literature's own repeated acknowledgment that
per-subject baselines are necessary
\cite{b6}\cite{b8}\cite{b9}\cite{b11}, even though, as
established in Section II, none of those systems quantify
that improvement themselves; the 6.0-percentage-point gain
reported in Table~\ref{tab_comparison} is, to the authors'
knowledge, among the first such quantified figures for this
sensor combination.

% ── SECTION 5 — DISCUSSION ───────────────────────────────
\section{Discussion}

\subsection{Personalization Impact on
Stress Detection}
The consistent improvement across every metric in
Table~\ref{tab_comparison}, and especially the 11.3
percentage-point gain in F1-Stress, is strong evidence that a
meaningful share of the static model's error in
Table~\ref{tab_cm} was attributable to inter-subject
variability rather than to any limitation of the feature set
or classifier itself. A population-level threshold necessarily
sits somewhere between individuals' true resting states; for
any given subject, that threshold is either too sensitive
(producing false Stress alarms during their normal resting
variability) or not sensitive enough (missing genuine stress
responses that, for that individual, are still well within the
population's "normal" range). Personalizing the reference point
that each new reading is compared against directly corrects
this mismatch, which is consistent with why the Stress class,
the class most directly affected by where the decision boundary
sits relative to an individual's true baseline, shows the
largest single improvement.

\subsection{EMA Continual Learning Behaviour}
Restricting the EMA update in Eq.~\ref{eq_ema} to
confidently-Normal windows only, and requiring a minimum of
three such windows per session, is a deliberate asymmetry: it
allows the baseline to track genuine long-term drift while
remaining structurally incapable of drifting \emph{toward} a
stressed state, since a session dominated by Stress
classifications simply contributes no update at all rather than
pulling the baseline toward a plausible-looking-but-wrong
"new normal." The choice of $\alpha = 0.2$ reflects a
deliberate stability-versus-responsiveness trade-off: a
larger $\alpha$ would adapt faster but would also let a single
atypical session (e.g. illness, unusually poor sleep) exert
outsized influence on the baseline, while a smaller $\alpha$
would better resist such outliers but track genuine
physiological change too slowly to be useful over realistic
usage timescales.

\subsection{Trend Monitoring as a Supplementary Signal}
The trend indicator (Eq.~\ref{eq_trend}) is deliberately scoped
as a minor addition on top of the personalization layers that
are this work's central contribution: it requires majority
agreement across monitored features before firing, which
suppresses false positives from ordinary single-feature noise,
but its real-world, session-level effectiveness remains to be
quantified on ring data (Section~IV).

\subsection{GSR Essentiality Confirmed}
The SWELL cross-dataset result (Table~\ref{tab_swell}) is one
of the more concrete findings in this work precisely because it
is a negative result under a fair comparison: the same
model-training procedure, restricted only to the features SWELL
actually provides, collapses from F1-macro 0.617 to 0.270. This
is not merely evidence that more features are better in
general; it specifically isolates electrodermal activity as
the feature domain responsible for the majority of this
model's discriminative power, which has a direct design
implication: any stress-detection wearable that omits EDA
sensing (as most consumer wrist-worn devices currently do, since
reliable EDA measurement requires two-point skin contact that
is awkward to achieve at the wrist) should be expected to
underperform substantially relative to a device, such as the
finger-worn ring described here, that retains it.

\subsection{Comparison with Related Work}
Table~\ref{tab_comparison2} situates this work against the
three most directly comparable systems identified in Section
II. None of the compared systems implement personalization or
continual adaptation, the central contributions of this work,
nor the supplementary trend indicator; this work's
combination of all three, on top of the same PPG+GSR+ACC
sensing modality used by these prior systems, represents the
first system in the reviewed literature to combine adaptive,
individualized modeling with a lightweight early-warning
supplement.

\begin{table}[htbp]
\caption{Comparison with Related Work}
\begin{center}
\begin{tabular}{|l|c|c|c|c|}
\hline
\textbf{Study} &
\textbf{Sensors} &
\textbf{Personal} &
\textbf{Adaptive} &
\textbf{Trend} \\
\hline
\cite{b10} &
PPG+GSR & No & No & No \\
\cite{b11} &
ECG+EDA & No & No & No \\
\cite{b6} &
Multi & No & No & No \\
\textbf{Ours} &
PPG+GSR+ACC &
\textbf{Yes} &
\textbf{Yes} &
\textbf{Yes} \\
\hline
\end{tabular}
\label{tab_comparison2}
\end{center}
\end{table}

\subsection{Limitations}
Four limitations of the present work should be stated plainly.
First, WESAD's stress condition is laboratory-induced via the
TSST protocol, which, while a validated and widely-used stress
induction method, is not identical to the more diffuse,
lower-intensity stress that characterizes everyday life;
generalization from TSST-induced stress to naturalistic daily
stress cannot be assumed without further validation. Second,
15 subjects is a modest sample for a personalization study in
particular, since the entire premise of the personalization
framework is that individual variability matters; a larger and more diverse
subject pool would strengthen confidence in the magnitude of
the personalization gains reported here. Third, the stress
trend detector's real-world effectiveness depends on
consistent sensor contact throughout wear; intermittent
electrode contact or sensor slippage on a worn device could
degrade the feature stability the trend detector relies on in
ways that do not occur in WESAD's controlled, securely-mounted
chest sensor placement. Fourth, the EMA continual-learning
layer requires a minimum of roughly three sessions with
sufficient Normal windows before it begins to meaningfully
diverge from the initial calibration, meaning its benefits are
not available to a user from their very first session.

% ── SECTION 6 — NOVEL CONTRIBUTIONS ─────────────────────
\section{Novel Contributions}
Table~\ref{tab_contributions} summarizes this work's seven
novel contributions, each mapped to the specific literature gap
identified in Section II (Table~\ref{tab_gaps}) that it
addresses. Collectively, these contributions span the full
pipeline: from how the ground-truth label itself is
constructed (N1, N6), through the personalized baseline
calibration and continual-adaptation mechanisms that form this
work's central contribution and are entirely new to this
domain (N2, N3), through a principled treatment of motion as
both a confound and a feature (N5), a lightweight supplementary
trend-monitoring indicator (N4), and the hardware form factor
that makes the electrodermal sensing shown to be essential in
Section~\ref{sec_results_obj1} practical to wear continuously
(N7).

\begin{table}[htbp]
\caption{Novel Contributions and
Literature Justification}
\begin{center}
\begin{tabular}{|c|p{3cm}|p{3cm}|}
\hline
\textbf{N} &
\textbf{Contribution} &
\textbf{Gap Addressed} \\
\hline
N1 & PSS-aligned stress labelling
of WESAD &
Only 11\% use PSS \cite{b15} \\
N2 & Personalized 2-min baseline
calibration &
Individual variability
unsolved \\
N3 & EMA continual learning
(Normal windows only) &
No adaptive system
exists \\
N4 & Feature-aware stress trend
detector (2.5 min lookback) &
No early warning
in literature \\
N5 & MPU-6050 dual role:
gating + feature &
Motion confound
unaddressed \\
N6 & LOSO cross-validation
standard &
Only 20\% use LOSO
\cite{b15} \\
N7 & Smart ring form factor
with WiFi/MQTT streaming &
Emerging direction
\cite{b6} \\
\hline
\end{tabular}
\label{tab_contributions}
\end{center}
\end{table}

% ── SECTION 7 — CONCLUSION ───────────────────────────────
\section{Conclusion}

This paper presented a wearable physiological stress-detection
system built around a rigorously LOSO-validated binary
Normal/Stress classifier and a personalization framework
(per-user baseline calibration combined with exponential-moving-
average continual learning) designed to address the
individual-variability and staleness gaps identified as common
across the fifteen studies reviewed in this work.

Personalization improved every reported metric over the static
baseline: F1-macro rose from 0.617 to 0.628, Ordinal Mean
Absolute Error fell from 0.357 to 0.194, overall accuracy
improved by 6 percentage points, and F1-Stress specifically
improved by 11.3 percentage points, the largest gain of any
single metric and direct evidence that individual baseline
normalization disproportionately helps precisely the class a
population-level model struggles with most. A lightweight
supplementary trend-monitoring indicator, layered on top of
this personalization framework, has so far been validated only
on controlled scenarios; its session-level effectiveness on
the deployed hardware remains near-term future work, described
below.

A cross-dataset check against SWELL, using only the HRV
features available in that dataset, showed a substantial
performance collapse relative to the full model (F1-macro 0.617
to 0.270), confirming that electrodermal activity is not a
convenient addition but a necessary input to this pipeline,
directly justifying the finger-worn ring form factor over a
wrist-only alternative that could not reliably sense it.

Beyond the IoT ring hardware validation reported in Section IV,
future work priorities include extending personal-baseline
adaptation to per-user decision thresholds (rather than only
feature normalization), a longitudinal study of the EMA layer's
behavior over months rather than the handful of sessions
evaluated here, and validation of the HRV pipeline against a
clinical-grade ECG reference to quantify any accuracy gap
introduced by the ring's PPG-based sensing relative to the
chest-ECG-based sensing used for training.

% ── ACKNOWLEDGMENT ───────────────────────────────────────
\section*{Acknowledgment}
The author thanks their supervisor for guidance throughout
this research, and the Sri Lanka Institute of Information
Technology, Department of Information Technology, for
supporting this work.

% ── REFERENCES ───────────────────────────────────────────
\begin{thebibliography}{00}

\bibitem{b1}
A. Taskasaplidis et al.,
``Review of stress detection methods using
wearable sensors,''
\textit{IEEE Access}, 2024.

\bibitem{b2}
K. M. Dalmeida and G. L. Masala,
``HRV features as viable physiological markers
for stress detection using wearable devices,''
\textit{Sensors}, vol. 21, 2021.

\bibitem{b3}
M. Haque et al.,
``State-of-the-art of stress prediction from
heart rate variability using AI,''
\textit{Cognitive Computation}, 2024.

\bibitem{b4}
I. Lazarou and T. Exarchos,
``Predicting stress levels using physiological
data,''
\textit{AIMS Neuroscience}, 2024.

\bibitem{b5}
M. Kyriakou et al.,
``Detecting moments of stress from measurements
of wearable physiological sensors,''
\textit{Sensors}, vol. 19, 2019.

\bibitem{b6}
R. Prasad and R. Upadhyaya,
``Stress monitoring using wearable technology,''
\textit{IEEE ICCA}, 2025.

\bibitem{b7}
Y. Mozafari et al.,
``Towards IoT-enabled multimodal mental stress
monitoring,''
\textit{IEEE COINS}, 2020.

\bibitem{b8}
M. V. Villarejo et al.,
``A stress sensor based on galvanic skin
response controlled by ZigBee,''
\textit{Sensors}, vol. 12, 2012.

\bibitem{b9}
H. Dybvik et al.,
``fNIRS, EEG, ECG, and GSR reveal physiological
stress,''
\textit{Frontiers in Human Neuroscience}, 2025.

\bibitem{b10}
D. Segono et al.,
``IoT-based stress monitoring using CNN for
HRV-GSR analysis,''
\textit{Sinkron Journal}, 2026.

\bibitem{b11}
R. Calvo et al.,
``Early detection of chronic stress using
wearable devices on WESAD,''
\textit{ICT4AWE}, 2025.

\bibitem{b12}
O. Nechyporenko et al.,
``Galvanic skin response and PPG for stress
recognition,''
\textit{Applied Sciences}, 2024.

\bibitem{b13}
S. Watharkar et al.,
``Real-time IoT-based physiological stress
assessment system,''
\textit{IRJET}, 2024.

\bibitem{b14}
S. Gedam et al.,
``A deep learning approach with wearable
physiological sensors to detect mental stress,''
\textit{Discover Artificial Intelligence}, 2026.

\bibitem{b15}
L. Pataca et al.,
``Use of machine learning for predicting stress
episodes based on wearable sensor data:
a systematic review,''
\textit{Computers in Biology and Medicine}, 2025.

\end{thebibliography}

\end{document}
